\chapter{Concrete objectives, lifting and rates}\label{comp:F}\label{F:chapter}
\renewcommand{\ModuleTag}{6}
These applications turn the common analytic mechanism into verifiable objective classes. The chapter also derives an affine-fiber selection formula and quantitative rates.


% SOURCE_BEGIN F.span.1
\section{Centered smooth-plus-sublinear objectives}

This chapter collects nonredundant structural consequences of the general theory. The finite-polyhedral-partition and powered-gauge radial arguments themselves are reproduced in Chapter~\ref{comp:R}. The discussion here gives the global, nonunique, explicit-basin and quantitative clauses, which go beyond a local uniqueness theorem.

\begin{theorem}[Centered smooth-plus-sublinear convergence]\label{F:centered}
Let $F(x)=\psi(x)+\gamma(x-x_*)$, where $\psi$ is convex and locally $C^{1,1}$ on an open convex neighborhood of a compact initial sublevel, and $\gamma$ is finite convex and positively homogeneous of degree one. Suppose $x_*$ minimizes $F$. Every infinite admissible full-BFGS run has $F(x_k)\to F(x_*)$. If $x_*$ is unique, then $x_k\to x_*$. A representation only near a unique minimizer gives a low-starting local theorem by Lemma~\ref{M:localization}.
\end{theorem}
\begin{proof}
Translate $x_*=0$ and let $R$ bound the initial sublevel. At differentiable endpoints, $\gamma$ is differentiable because $F$ and $\psi$ are. Write $y=y_\psi+y_\gamma$, and set
\[
 a=y_\psi^Ts,\quad b=y_\gamma^Ts,\quad
 \xi=y_\psi^Tx',\quad\zeta=y_\gamma^Tx'.
\]
Then $a,b\ge0$ and $t=a+b$. Convex local $L$-smoothness gives $\|y_\psi\|^2\le La$, hence $\xi^2\le LR^2a$. Here $L$ is a uniform constant chosen for the compact endpoint set inside the open convex neighborhood. A local proof subtracts a supporting tangent plane and applies the descent estimate at a short step opposite the gradient difference, with the test step kept inside a larger buffered neighborhood. Applying the resulting Bregman estimate in both directions gives the stated inequality. The constant may be enlarged to keep those test steps inside the neighborhood.

If $h\in\partial\gamma(x)$, homogeneity and the supporting inequality give $h^Tx=\gamma(x)$ and $h^Tz\le\gamma(z)$. Hence $0\le\zeta\le b$. All subgradients of $\gamma$ are bounded by some $G_\gamma$, so $\zeta\le2G_\gamma R$. Weighted Cauchy--Schwarz gives
\[
 \frac{(y^Tx')^2}{t}
 \le\frac{\xi^2}{a}+\frac{\zeta^2}{b}
 \le LR^2+2G_\gamma R.
\]
Terms with zero denominator are zero since their numerators vanish. Apply the fixed-reference energy bound and the reciprocal-average theorem in Chapter~\ref{comp:M}. Compactness and uniqueness give the iterate conclusion.
\end{proof}
The original summands need not separately be minimized at the center. Optimality implies $-\nabla\psi(x_*)\in\partial\gamma(0)$. After translating, subtract the smooth tangent plane from $\psi$ and add it to $\gamma$; both adjusted summands are then nonnegative and minimized at zero. The objective itself is unchanged by this rewriting.

% SOURCE_END F.span.1


% SOURCE_BEGIN F.span.2
\section{Smooth increasing transforms of coercive gauges}
\begin{theorem}[Transformed-gauge convergence]\label{F:transform}
Let $\gamma$ be finite convex, degree-one homogeneous and coercive. Let $\Phi:[0,\infty)\to\mathbb R$ be continuous, convex, strictly increasing, with $\Phi(0)=0$, $\Phi\in C^2((0,\infty))$ and $\Phi'>0$ there. Then every infinite admissible full-BFGS run on $F(x)=\Phi(\gamma(x-x_*))$ converges to $x_*$.
\end{theorem}
\begin{proof}
Coercivity of $\gamma$ and unboundedness of $\Phi$ give compact sublevels. Translate $x_*=0$. Suppose $F_k\downarrow L>0$. Euler's identity and the chain rule give
\[
 a_k=g_k^Tx_k=\gamma(x_k)\Phi'(\gamma(x_k))=\Theta(F_k),
\quad\Theta(t)=\Phi^{-1}(t)\Phi'(\Phi^{-1}(t)).
\]
On $[L,F_0]$, the derivative
\[
 \Theta'(t)=1+\frac{r\Phi''(r)}{\Phi'(r)},\qquad r=\Phi^{-1}(t),
\]
is bounded. Thus $|a_{k+1}-a_k|\le C\Delta_k$. Since
\[
 y_k^Tx_{k+1}=d_k+a_{k+1}-a_k,
\]
\eqref{M:basic} implies
\[
 R_k\le\frac{(1+C)^2}{(1-c_2)c_1}\Delta_k.
\]
The corrections are summable, contradicting a positive gap in \eqref{M:energybound}. Hence $F_k\to0$, and coercivity and strict increase give $x_k\to0$.
\end{proof}
The same correction estimate can be obtained without introducing an inverse: on a compact positive gauge annulus, bound the derivative of $r\Phi'(r)$ and bound $\Phi'$ away from zero, then apply the scalar mean-value theorem twice. The formal convergence argument uses this annulus route; the inverse identities and derivative formula displayed above are separately proved. Theorem~\ref{F:transform} includes cases whose scalar second derivative is unbounded at zero. Finite sums of differently powered gauges are instead covered by the radial argument in Chapter~\ref{comp:R}, including $1<p<2$.

% SOURCE_END F.span.2


% SOURCE_BEGIN F.span.3
\section{Finite max-affine terms and a certified basin}
Let $F=\psi+p$, with $\psi$ convex locally $C^{1,1}$ and $p(x)=\max_i\ell_i(x)$, $\ell_i(x)=a_i^Tx+b_i$, a finite real-valued maximum. At a unique minimizer $x_*$, put $I_*=\{i:\ell_i(x_*)=p(x_*)\}$. Inactive pieces have positive gaps, so on a sufficiently small ball
\begin{equation}\label{F:active}
 p(x_*+z)=p(x_*)+\max_{i\in I_*}a_i^Tz.
\end{equation}
Indeed, with any fixed $i_0\in I_*$, the difference between its value and an inactive piece $j$ is
\[
 \delta_j+(a_{i_0}-a_j)^Tz,\qquad \delta_j=p(x_*)-\ell_j(x_*)>0,
\]
which stays positive if $\|z\|\le\delta_j/(2\|a_{i_0}-a_j\|)$; a parallel inactive piece stays dominated everywhere. This exact local identity and Theorem~\ref{F:centered} prove local convergence without strict complementarity or Hessian nondegeneracy.

For an explicit sufficient level, define the closed active region
\[
 P_j=\{x:\ell_j(x)=p(x)\}
 =\{x:(a_j-a_i)^Tx\ge b_i-b_j\ \forall i\}
\]
and
\begin{equation}\label{F:basin}
 \tau_* =\min_{j\notin I_*}\inf_{x\in P_j}F(x),
\end{equation}
with empty infima or minima equal to $+\infty$. Then $\tau_*>F_*$. To prove it, note that unique-minimizer convexity gives compact sublevels; a finite infimum on a nonempty closed $P_j$ is attained. It cannot equal $F_*$ because $x_*\notin P_j$. Finitely many inactive pieces preserve the strict gap.

Every admissible run with $F(x_0)<\tau_*$ converges. On its full compact initial sublevel, only indices in $I_*$ can be active, so \eqref{F:active} and the centered correction estimate apply throughout. A tail that first enters this sublevel enjoys the same conclusion for its current arbitrary SPD matrix. Each inner infimum is a convex program over the displayed linear inequalities.

% SOURCE_END F.span.3


% SOURCE_BEGIN F.span.4
\section{Residual penalties and affine-fiber lifting}

For a scalar affine residual nonzero at the solution, its absolute-value or hinge term is locally affine. A zero optimal residual produces a centered sublinear term. The same decomposition applies to an affine-composed Euclidean norm: a nonzero optimal vector residual is locally smooth, and a zero residual becomes a centered norm. Finite sums with a convex locally $C^{1,1}$ smooth part reduce to Theorem~\ref{F:centered}. For a more general gauge, local $C^{1,1}$ smoothness at that particular nonzero residual must be checked; a spectral penalty may fail it.

\begin{theorem}[Exact linear-image lifting]\label{F:lifting}
Let $A\in\mathbb R^{r\times n}$ have full row rank and $F(x)=\bar F(Ax)$. Suppose the projected admissible run converges to a unique optimizer $z_*$ of $\bar F$, under any of the retained sufficient hypotheses. Then
\[
 x_k\to x_0+H_0A^T(AH_0A^T)^{-1}(z_*-Ax_0).
\]
\end{theorem}
\begin{proof}
With $z_k=Ax_k$ and $\bar H_k=AH_kA^T$, projection of the inverse-BFGS update gives exactly full BFGS on $\bar F$, with the same step lengths and line-search comparisons. Let $S=\operatorname{range}A^T$. Since gradients and secant differences lie in $S$, the update gives $H_{k+1}S\subset H_kS$; both spaces have dimension $r$, so equality holds. All displacements lie in $H_0S$. The restriction of $A$ to that space has right inverse
\[
 L_0=H_0A^T(AH_0A^T)^{-1}.
\]
Thus $x_k-x_0=L_0(z_k-z_0)$, proving the result. This is the $H_0^{-1}$-metric projection onto the affine fiber.
\end{proof}

% SOURCE_END F.span.4


% SOURCE_BEGIN F.span.5
\section{Finite-horizon rates and strengthened asymptotic corollaries}

Let $S_N=\sum_{k<N}q_k^{-1}$. Retain the explicit lower bound \eqref{M:finitehorizon}, which is meaningful at a fixed horizon. All rates below concern fixed dimension, problem, initialization and line-search constants, and count accepted iterations only. Big-$O$ versions of these rates appeared in the source notes. The little-$o$ versions below follow from Theorem~\ref{M:superquadratic}.

\paragraph{Uniformly bounded reference correction.}
If $R_k\le\bar R$, then \eqref{M:energybound} gives
\[
 F(x_N)-F_*\le\sqrt{\frac{E_0+\bar RN}{S_N}},
\]
and along runs with bounded gradients this is $o(N^{-1/2})$, because $S_N/N^2\to\infty$. The same estimate holds for a linear cumulative correction budget.

\paragraph{Bounded-below degree-one objectives.}
Chapter~\ref{comp:H} proves
\[
 \sum_{k<N}\frac{F(x_k)^2}{q_k}\le\Phi_0,
\quad \Phi_0=x_0^TH_0^{-1}x_0+\frac{1-c_1}{c_1}F(x_0).
\]
Hence $F(x_N)\le\sqrt{\Phi_0/S_N}=o(N^{-1})$. Coercivity gives $F(x)\ge\mu\|x\|$ and the same order for the iterate norm. The little-$o$ conclusion uses the superquadratic theorem; the finite-horizon bound retains its explicit constants.

\paragraph{Quadratic plus coercive gauge.}
For $F(u,v)=u^TQu/2+\gamma(v)$, $Q\succ0$, let $\mu\|v\|\le\gamma(v)\le\Lambda\|v\|$. Split $t=a+b$ with
\[
 a=\delta u^TQ\delta u,\quad b=(h'-h)^T\delta v,\quad
 \xi=(Q\delta u)^Tu',\quad\zeta=(h'-h)^Tv',
\]
where $h\in\partial\gamma(v)$ and $h'\in\partial\gamma(v')$. Then
\[
 \xi^2/a\le u'^TQu',\qquad0\le\zeta\le b,
 \qquad\zeta\le(1+\Lambda/\mu)\gamma(v').
\]
Thus
\[
 R_k\le C_\gamma F(x_{k+1}),\qquad
 C_\gamma=\max\{2,1+\Lambda/\mu\}.
\]
Choose $a_0>0,N_0$ with $S_N\ge a_0N^2$ for $N\ge N_0$, and put
\[
 A_N=E_0+C_\gamma\sum_{j=1}^NF(x_j).
\]
Then $F(x_N)\le\sqrt{A_N}/(\sqrt{a_0}N)$, and
\[
 \sqrt{A_N}-\sqrt{A_{N-1}}
 =\frac{C_\gamma F(x_N)}{\sqrt{A_N}+\sqrt{A_{N-1}}}
 \le\frac{C_\gamma}{\sqrt{a_0}N}.
\]
Summing first gives the bounds $\sqrt{A_N}=O(\log N)$ and $F(x_N)=O(\log N/N)$. For the strengthened conclusion, keep the exact denominator:
\[
 \sqrt{A_N}-\sqrt{A_{N-1}}\le \frac{C_\gamma}{\sqrt{S_N}}=o(N^{-1}).
\]
A sum of $o(1/N)$ increments is $o(\log N)$: for any fixed $\epsilon>0$, bound the sufficiently late increments by $\epsilon/N$, then divide their sum by $\log N$ and let $\epsilon\downarrow0$. Thus $\sqrt{A_N}=o(\log N)$. Dividing by $\sqrt{S_N}=N\sqrt{S_N/N^2}$ yields
\[
 \boxed{F(x_N)=o(\log N/N).}
\]
These strengthened asymptotic corollaries follow from already recorded estimates.

% SOURCE_END F.span.5


% SOURCE_BEGIN F.span.6
\section{Matrix limits and line-search existence are different questions}

In the following runs, the matrix limit depends on the initial matrix. For $u^2+|v|$, choose $A>0$, $v_0>0$ and $0<\rho<(1-c_1)/(1+c_1)$. With
\[
 x_0=(0,v_0),\quad
 H_0=\operatorname{diag}\left(A,\frac{1+\rho}{2\rho}v_0\right),\quad\alpha_k=2\rho,
\]
ordinary BFGS gives
\[
 x_k=(0,(-\rho)^kv_0),\quad
 H_k=\operatorname{diag}\left(A,\frac{1+\rho}{2\rho}\rho^kv_0\right).
\]
The Armijo ratio is $(1-\rho)/(1+\rho)>c_1$ and the weak-Wolfe ratio is $-1$. The unexcited tangential entry tends to its arbitrary initial value $A$, not necessarily the inverse true tangential curvature.

For a finite convex line restriction $\varphi$ bounded below on the positive ray, suppose $\varphi'(0)=a<0$. Choose $\sigma=(c_1+c_2)/2$ and set $\psi(t)=\varphi(t)-\sigma at$. Its derivative at zero is $(1-\sigma)a<0$, while the lower bound on $\varphi$ makes $\psi(t)\to+\infty$ as $t\to+\infty$. Choose a compact interval $[0,M]$ whose interior contains a minimizer $z>0$ of $\psi$, with $\psi(z)<\psi(0)$. Then
\[
 \varphi(z)<\varphi(0)+\sigma az<\varphi(0)+c_1az.
\]
Continuity preserves strict Armijo on an interval immediately to the right of $z$. For each $t$ in a sufficiently small such interval, minimality and convexity give
\[
 \varphi'(t)\ge\frac{\varphi(t)-\varphi(z)}{t-z}\ge\sigma a>c_2a
\]
whenever the scalar derivative exists. Thus there is a nonempty open interval of strict Armijo points whose scalar derivatives satisfy strict weak Wolfe. If full-space nondifferentiability occurs at only finitely many positive ray times, choose a point in that interval outside this finite set; its full gradient gives a differentiable acceptable step. This is the tilted-minimum construction. It applies to finite affine absolute-value/hinge sums and suitable Euclidean residual sums through their proved full-space bad-ray-set finiteness.

For arbitrary max-functions there is an exception. The coercive function
\[
 h(x,y,z)=\max\{-x,y+z,-y+z,x/4,-z\}
\]
has unique minimizer zero. At $(-2,0,1)$ with $H_0=I$, the direction is $(1,0,0)$ and the ray function is $\max\{2-\alpha,1,(\alpha-2)/4\}$. For $(c_1,c_2)=(1/4,1/2)$, times $0<\alpha<1$ fail weak Wolfe; $1\le\alpha\le6$ are full-space nondifferentiable; and $\alpha>6$ fails Armijo, which would require $\alpha\le5$. Scaling the point and matrix places the same finite obstruction arbitrarily close to the minimizer. It shows why the positive theorems quantify over infinite admissible runs: no such run starts from these data.
% SOURCE_END F.span.6


\section*{Relation to earlier work}
Guo--Lewis~\cite[Theorem~4.1 and Corollary~4.4]{GL18} uses smooth extension and proves a weaker subsequential conclusion for $u^2+|v|$. The theorem in this chapter gives full convergence for that example within its centered family. Its smooth secant estimate uses established cocoercivity~\cite[Theorem~2.5]{WW}: local regularity is used on an open convex neighborhood, with compact buffers and, where needed, an enlarged constant. It does not require a globally Lipschitz gradient or minimization of each original summand at the center.

Invertible affine invariance is established in Lewis--Overton~\cite[Section~2]{LO} and Xie--W\"achter~\cite[Theorem~II.1]{XW}. The rectangular reduction in this chapter additionally identifies the point selected on an affine fiber. Sch\"opfer's smooth noncompact convergence results~\cite[Theorems 3.4 and 3.7]{S16} do not identify that selector.

Related quantitative work includes smooth strongly convex Armijo--Wolfe rates~\cite{FND_JJM24AW}, smooth exact-search rates~\cite{FND_JJM26EX}, and smooth self-concordant analysis~\cite{FND_JM25}. These assumptions and line searches differ from the nonsmooth settings in this chapter. Their matrix-potential techniques are precedents, but their rate theorems do not imply these estimates.

Lewis--Overton~\cite[Theorem~4.5 and Proposition~4.11]{LO} supplies scalar Armijo--weak-Wolfe existence theory. Differentiability of the restriction to a ray does not imply full-space differentiability at acceptance. The existence argument and the finite obstruction in this chapter preserve that distinction.
